\section{Exact algebraic and semialgebraic ReLU encoding}
\label{sec:relu}

The graph of a ReLU is basic semialgebraic and has a particularly small
complementarity description.

\begin{lemma}[Exact ReLU complementarity]
\label{lem:relu-complementarity}
For \(z,a\in\R\),
\begin{equation}
  a=\ReLU(z)
  \quad\Longleftrightarrow\quad
  a\ge0,\quad a-z\ge0,\quad a(a-z)=0.
  \label{eq:relu-complementarity}
\end{equation}
\end{lemma}

\begin{proof}
Suppose first that \(a=\max\{0,z\}\).  If \(z\le0\), then \(a=0\), so
\(a\ge0\), \(a-z=-z\ge0\), and the product vanishes.  If \(z\ge0\), then
\(a=z\), so both weak inequalities hold and \(a-z=0\).

Conversely, the product equation implies either \(a=0\) or \(a-z=0\).  In the
first case, \(a-z\ge0\) implies \(z\le0\), whence
\(a=0=\max\{0,z\}\).  In the second case, \(a=z\) and \(a\ge0\), whence
\(a=z=\max\{0,z\}\).
\end{proof}

The three regimes are all present:
\begin{align}
  z<0 &\Longrightarrow a=0,\\
  z=0 &\Longrightarrow a=0,\\
  z>0 &\Longrightarrow a=z.
\end{align}
No disequality \(z\ne0\) is introduced.

\begin{lemma}[Equality-only polynomialization]
\label{lem:relu-polynomialization}
For \(z,a\in\R\),
\begin{equation}
 a=\ReLU(z)
 \quad\Longleftrightarrow\quad
 \exists u,v\in\R:
 a=u^2,\quad a-z=v^2,\quad a(a-z)=0.
 \label{eq:relu-polynomialization}
\end{equation}
\end{lemma}

\begin{proof}
By \cref{lem:relu-complementarity}, a ReLU assignment makes \(a\) and \(a-z\)
nonnegative.  Each is therefore a square in \(\R\).  Conversely, the two square
equalities imply the weak inequalities, and the complementarity equation then
allows an application of \cref{lem:relu-complementarity}.
\end{proof}

\begin{remark}[Projection is part of the semantics]
The equality variety in variables \((z,a,u,v)\) is not literally the graph in
the \((z,a)\)-plane.  Its projection to \((z,a)\) is the graph.  The slack
variables therefore belong to the existential local state.
\end{remark}

\begin{proposition}[Membership in \(\ExistsR\)]
\label{prop:nntp-in-existsr}
The language \(\ExactNNTP\) belongs to \(\ExistsR\).
\end{proposition}

\begin{proof}
Introduce real variables for the shared parameters, every sample-local
preactivation and activation, and the two square slacks of each ReLU.  Affine
transitions with trainable weights are polynomial equalities of degree at most
two in the unknowns.  Exact targets are polynomial equalities.  Replace every
ReLU by \cref{eq:relu-polynomialization}.  The number of variables and
constraints and the coefficient encoding length grow polynomially with the
explicit instance encoding.  The resulting existential conjunction of
polynomial equalities is true exactly when the instance is feasible.
\end{proof}

\subsection{The switching boundary and singularity}

The complementarity polynomial
\begin{equation}
  g(z,a)=a(a-z)
\end{equation}
has gradient
\begin{equation}
  \nabla g(z,a)=(-a,,2a-z).
\end{equation}
At \((z,a)=(0,0)\), this gradient vanishes.  Thus a Jacobian-based regularity
argument cannot cover the switching point merely because the global K\"ahler
module exists there.  K\"ahler differentials are defined at singular points,
but their fibers can change rank and acquire information not present on a
smooth stratum.  The exact semantics above remains valid at the boundary; only
the first-order regularity changes.
